\documentclass[11pt]{article}
\usepackage{fontspec,amsmath,amssymb,amsthm,hyperref}
\usepackage[margin=1in]{geometry}
\setlength{\emergencystretch}{3em}
\title{Prove Gibbs stationarity by refining invariant Metropolis chains}
\author{LangevinLean mathematical exposition}
\date{}
\begin{document}\maketitle
\noindent Source module: \texttt{UnadjustedLangevin/GibbsStationarity.lean}. Compile this proof document with XeLaTeX or LuaLaTeX. The accompanying HTML places the complete Lean declarations beside every proof.


A Brownian Langevin integral solution started from Gibbs stays Gibbs. The proof uses convergent Euler approximations and exactly invariant MALA chains, with no weak-PDE uniqueness assumption.


\section{Identify L\ensuremath{{}^{2}} Euler endpoints with the actual ULA laws}


Induction identifies the representative of the lifted chain with the pointwise noise-driven recursion. The Markov-law theorem then gives \(\operatorname{Law}(X_n)=P_h^n\operatorname{Law}(X_0)\) when the innovations are iid standard Gaussian vectors and independent of the initial state. If the additive noise satisfies \(B((k+1)h)-B(kh)=\sqrt{2h}Z_k\), induction also identifies the additive Euler chain with this lifted chain.



\noindent Lean declarations: \texttt{hilbert\_lpChain\_coe}, \texttt{hilbert\_lpChain\_law\_eq\_ulaLawStep}, \texttt{additiveEulerChain\_eq\_hilbert\_lpChain}.

\section{Use normalized Brownian increments}


Brownian additive increments telescope across the grid and equal \(\sqrt{2h}\) times the normalized Euler innovations. Therefore additive Euler endpoints driven by Brownian noise are the lifted ULA endpoints. The scalar coordinate normal laws and the stated joint coordinate-time independence imply iid standard Gaussian vector laws, using GaussianInnovationVectors. Independence from the initial state is retained explicitly.



\noindent Lean declarations: \texttt{brownianLangevinNoise\_grid\_increment}, \texttt{brownian\_additiveEulerChain\_eq\_lpChain}, \texttt{brownianEulerNoise\_vectorGaussian}.

\section{Pass the invariant Gibbs marginal to the continuous limit}


Fix \(T>0\) and choose \(h_n=T/(n+1)\). The Euler theorem gives \(X_{n+1}^{h_n}\to Y_T\) in \(L^2\). Couple each Euler endpoint, started at Gibbs, to the MALA chain at the same mesh. Exact MALA invariance makes the second marginal Gibbs; the rejected-displacement bound gives

\[\mathbb E\|X_{n+1}^{h_n}-M_{n+1}^{h_n}\|\le(n+1)Ch_n\sqrt{h_n}=CT\sqrt{h_n}\longrightarrow0.\]

For all sufficiently large \(n\), the mesh satisfies both \(h_n\le1\) and the stable-step restriction. Characteristic-function comparison between these couplings and continuity under the \(L^2\) limit identifies \(\operatorname{Law}(Y_T)=\pi\). At \(T=0\), the integral equation directly gives \(Y_0=X_0\). Strong convexity supplies Gibbs moments and the finite rejection constant; the theorem also displays Brownian increment independence for every positive mesh.



\noindent Lean declarations: \texttt{lp\_langevin\_gibbs\_stationary}.

\end{document}
